\documentclass[12pt]{article}
\usepackage[T1]{fontenc}
\usepackage[utf8]{inputenc}
\usepackage{lmodern}
\usepackage{textcomp}
\usepackage{enumitem}
\usepackage{geometry}
\usepackage{graphicx}
\usepackage{amsmath, amssymb}
\geometry{margin=1in}
\begin{document}
\begin{center}
Sociology - Undergraduate Level Assessment 25 \
Total Marks: 40 \
Answer all questions.
\end{center}

\section*{Multiple Choice (10 marks)}
\begin{enumerate}[label=\arabic*.]
\item It is difficult to ascertain the true extent of domestic violence because:
\begin{enumerate}[label=(\alph*)]
    \item there is a large 'dark figure' of unreported incidents
    \item the changing definitions of legal categories have made it harder to convict offenders
    \item researchers are not allowed access to official statistics
    \item there is no valid or reliable way of researching such a sensitive topic
\end{enumerate}
\item In contemporary societies, social institutions are:
\begin{enumerate}[label=(\alph*)]
    \item highly specialized, interrelated sets of social practices
    \item disorganized social relations in a postmodern world
    \item virtual communities in cyberspace
    \item no longer relevant to sociology
\end{enumerate}
\item Pilcher (1999) identified soap operas as a 'feminine genre' of media because:
\begin{enumerate}[label=(\alph*)]
    \item most of the characters in soap operas are women
    \item they represent images of women as both domesticated and independent
    \item they alienate women and appeal to an audience of men
    \item female television producers are most likely to work in this area
\end{enumerate}
\item Theories of racialized discourse suggest that:
\begin{enumerate}[label=(\alph*)]
    \item race is an objective way of categorizing people on biological grounds
    \item the idea of race is socially constructed through powerful ideologies
    \item race relations in Britain and America can be traced back to colonial times
    \item people choose their racial identity and this becomes fixed
\end{enumerate}
\item The term 'collective consumption' (Castells 1977) refers to:
\begin{enumerate}[label=(\alph*)]
    \item the privatization of public services by the Conservative government
    \item the lifestyle practice of shopping in peer groups
    \item the form of tuberculosis suffered by those who collect stamps
    \item the provision of health, housing, and education services by the state
\end{enumerate}
\end{enumerate}

\section*{True / False (10 marks)}
\textit{Answer True or False}
\begin{enumerate}[label=\arabic*.]
\setcounter{enumi}{5}
\item True or False: The 'correspondence principle' posits that schools prepare children for work by instilling obedience and conformity.
\item True or False: Telework primarily involves operating a personal business through telephone helplines or hotlines.
\item True or False: A population pyramid with a slightly larger birth rate than death rate will have parallel sides.
\item True or False: Scott (1991) defined the 'power elite' as a state elite primarily composed of members from a dominant power bloc.
\item True or False: Cultural restructuring focuses solely on preserving historical landscapes without any economic or advertising initiatives.
\end{enumerate}

\section*{Long Form Response (20 marks)}
\textit{Answer each question with a clear and complete explanation.}
\begin{enumerate}[label=\arabic*.]
\setcounter{enumi}{10}
\item Discuss the implications of Sullivan's (2000) study on gender roles in housework, particularly focusing on the conditions under which men engage more in household labor.
\item Discuss the key processes of decision-making in the USA as identified by Domhoff, and critically analyze why the exploitation process is not included in his framework.
\end{enumerate}

\vfill
\noindent\textit{End of Paper}
\end{document}